% ============================================================================
% PLANTILLA — CLASE (diapositivas de sesión)
% ----------------------------------------------------------------------------
% Presentación de una clase. SIN diseño: todo lo pone academic-beamer + el tema
% «Academic». Solo se elige la clase y se rellena el contenido.
% Compilar:  scripts/build.sh clase.tex   (XeLaTeX)
% ============================================================================
\documentclass{academic-beamer}

\title{Título de la clase}
\subtitle{Unidad · tema}
\author{Edison Achalma B.Sc. Econ.}
\institute{Universidad Nacional de San Cristóbal de Huamanga}
\date{Semestre 2026-I}

\begin{document}

\begin{frame}\titlepage\end{frame}

\begin{frame}{Objetivos de la sesión}
\begin{itemize}
  \item Objetivo de aprendizaje 1
  \item Objetivo de aprendizaje 2
  \item Objetivo de aprendizaje 3
\end{itemize}
\end{frame}

\section{Primer tema}
\begin{frame}{Idea principal}
\begin{itemize}
  \item Punto clave
  \begin{itemize}\item Detalle o ejemplo\end{itemize}
\end{itemize}
\begin{block}{Definición}
Enunciado formal del concepto.
\end{block}
\end{frame}

\begin{frame}{Formalización}
Ecuación o modelo: $\displaystyle y = f(x;\theta)$.
\begin{nota} Observación o matiz importante para la clase. \end{nota}
\end{frame}

\begin{frame}[fragile]{Aplicación en software}
\begin{codigo}[language=R]
modelo <- lm(y ~ x, data = datos)
summary(modelo)
\end{codigo}
\end{frame}

\begin{frame}{Síntesis}
\begin{itemize}
  \item Idea de cierre 1
  \item Idea de cierre 2
\end{itemize}
\end{frame}

\end{document}
